\FloatBarrier
\noindent Solar radiation pressure (SRP) provides a continuous, propellant-free means of modifying spacecraft orbits, and has been demonstrated in missions including IKAROS \citep{tsuda2011ikaros,sawadaMissionReportSolar2011}, LightSail~2 \citep{spencerLightSail2Solar2021}, and the Advanced Composite Solar Sail System \citep{wilkieOverviewNASAAdvanced2021}. Recent reviews summarise the current state of solar sailing technology and mission heritage \citep{johnsonSolarSailsTechnology2012,yakezerashviliRecentAdvancesSpace2025}. While active solar sailing enables controlled trajectory shaping, the possibility of achieving meaningful orbital energy change under purely passive configurations remains less well understood, particularly in low Earth orbit (LEO) where atmospheric drag and gravitational perturbations are significant. The use of solar sails in LEO has been explored for deorbiting \citep{stohlmanDevelopmentDeorbitsailFlight2014,lucking2012solar} and for altitude maintenance \citep{berthetSunflowerMissionConcept2023}, but the short-horizon interaction between eclipse timing and SRP work balance has received limited attention.

\vspace{0.5em}
\par\noindent Prior work has examined the coupled effects of solar radiation pressure and atmospheric drag on high area-to-mass spacecraft. In particular, \citet{colombomcinnes} demonstrated that asymmetric SRP, including the influence of Earth eclipse, can be exploited to balance drag losses and produce long-lived Earth-centred orbits without active control. The influence of SRP on high area-to-mass ratio objects has also been studied in the context of long-term orbit evolution and resonance structures \citep{valk2008semi}. These analyses highlight the importance of SRP–drag interaction and shadow-induced asymmetries in shaping orbital evolution. However, much of the existing SRP–drag literature emphasizes long‑term or averaged behavior and equilibrium‑like manifolds or ``balance'' conditions, rather than isolating short‑horizon regimes in which the timing of eclipse arcs alters the sign and magnitude of orbit‑averaged SRP work.

\vspace{0.5em}
\par\noindent In near-polar LEO, the interaction between SRP, atmospheric drag, Earth’s oblateness ($J_2$), and eclipse timing produces a dynamically coupled environment in which small perturbations can lead to non-intuitive orbital behaviour over short timescales \citep{prieto2014drag}. In particular, a passively oriented, spin-stabilised sail experiences alternating phases of positive and negative SRP work over each orbit, which would be expected to largely cancel in the absence of asymmetries \citep{mcinnes1999}.

\vspace{0.5em}
\par\noindent However, eclipse introduces a temporal modulation of SRP by removing solar forcing over specific orbital arcs \citep{hubaux2013shadow}. Depending on the orbital geometry, this can selectively suppress phases in which SRP performs negative work, altering the orbit-averaged energy balance. This raises the question of whether a net increase in semi-major axis can arise over a bounded time horizon under restricted conditions, without active attitude control.

\vspace{0.5em}
\par\noindent This work investigates that question by isolating and characterising a transient, geometry-dependent energy-pumping regime in near-polar LEO for spin-stabilised solar sails. The aim is not to propose a mission architecture or demonstrate sustained orbit raising, but to identify where the regime exists, explain the eclipse-modulated SRP work-imbalance mechanism, and quantify the conditions under which it appears and breaks down. Using a reduced-order dynamical model and parameter-space sweeps over local time of descending node (LTDN) and sail orientation, we show that short-horizon semi-major-axis growth can occur within a narrow alignment envelope, but that the effect is inherently self-limiting through coupled eccentricity growth, increased drag exposure near perigee, and gradual drift in Sun-relative geometry.
\FloatBarrier